\documentclass{article}
\usepackage[utf8]{inputenc}
\usepackage{amsmath}
\usepackage{geometry}
\usepackage{fancyhdr} % Paket für Kopf- und Fußzeilen
\geometry{a4paper, margin=1in}

\pagestyle{fancy}
\fancyhf{}
\rhead{Lampert Mathias}
\lhead{Ausgewähle Kapitel der Elektronik}
\cfoot{\thepage}

\begin{document}

\section*{Detaillierte Herleitung der Schaltverlustenergie bei ohmscher Last}

\subsection*{1. Definition der Verlustleistung}
Die momentane Verlustleistung $p(t)$ an einem Halbleiterschalter (z. B. MOSFET) während des Schaltvorgangs ist definiert als das Produkt aus der über dem Bauteil abfallenden Spannung $u_{DS}(t)$ und dem hindurchfließenden Drainstrom $i_D(t)$:
\begin{equation}
p(t) = u_{DS}(t) \cdot i_D(t)
\end{equation}

\subsection*{2. Modellierung des Schaltvorgangs}
Für eine rein ohmsche Last wird angenommen, dass der Strom und die Spannung während der Schaltzeit $t_s$ linear verlaufen. Wir betrachten den Einschaltvorgang im Zeitintervall $t \in [0, t_s]$:

\begin{itemize}
    \item \textbf{Linearer Stromanstieg}: Der Strom steigt von $0$ auf den Maximalwert $I_D$ an:
    \begin{equation}
    i_D(t) = I_D \cdot \frac{t}{t_s}
    \end{equation}
    
    \item \textbf{Linearer Spannungsabfall}: Die Spannung sinkt von der Versorgungsspannung $U_b$ auf $0$ ab:
    \begin{equation}
    u_{DS}(t) = U_b \cdot \left( 1 - \frac{t}{t_s} \right)
    \end{equation}
\end{itemize}



\subsection*{3. Aufstellen des Integrals}
Die gesamte während eines Schaltvorgangs in Wärme umgesetzte Energie $W_v$ berechnet sich durch Integration der Momentanleistung über die Zeitdauer $t_s$:
\begin{equation}
W_v = \int_{0}^{t_s} p(t) \, dt = \int_{0}^{t_s} \left[ U_b \cdot \left( 1 - \frac{t}{t_s} \right) \right] \cdot \left[ I_D \cdot \frac{t}{t_s} \right] \, dt
\end{equation}

\begin{equation}
W_v = U_b \cdot I_D \int_{0}^{t_s} \left( 1 - \frac{t}{t_s} \right) \cdot \frac{t}{t_s} \, dt
\end{equation}

\begin{equation}
W_v = U_b \cdot I_D \int_{0}^{t_s} \left( \frac{t}{t_s} - \frac{t^2}{t_s^2} \right) \, dt
\end{equation}

\subsection*{4. Integration und Auswertung}
Wir bestimmen die Stammfunktion bezüglich der Variablen $t$:
\begin{equation}
W_v = U_b \cdot I_D \cdot \left[ \frac{1}{t_s} \cdot \frac{t^2}{2} - \frac{1}{t_s^2} \cdot \frac{t^3}{3} \right]_{0}^{t_s}
\end{equation}

\begin{equation}
W_v = U_b \cdot I_D \cdot \left( \frac{t_s^2}{2 t_s} - \frac{t_s^3}{3 t_s^2} \right)
\end{equation}

\begin{equation}
W_v = U_b \cdot I_D \cdot \left( \frac{t_s}{2} - \frac{t_s}{3} \right)
\end{equation}

\begin{equation}
W_v = U_b \cdot I_D \cdot \left( \frac{3 t_s}{6} - \frac{2 t_s}{6} \right) = U_b \cdot I_D \cdot \frac{t_s}{6}
\end{equation}

\subsection*{5. Endergebnis}
Die Energie pro Schaltvorgang bei ohmscher Last beträgt:
\begin{equation}
W_v = \frac{1}{6} \cdot U_b \cdot I_D \cdot t_s
\end{equation}

Die mittlere Schaltverlustleistung $P_v$ bei einer Schaltfrequenz $f$ ergibt sich schließlich zu:
\begin{equation}
P_{v, Schalten} = W_v \cdot f = \frac{U_b \cdot I_D \cdot t_s \cdot f}{6}
\end{equation}

\end{document}